\myaddRA{\subsection{Instantiation I: Counter-State DPCM for \texorpdfstring{$r=1$}{r=1}}}
\label{sec:r1}

Under Theorem~\ref{thm:dpcm-states}, the sufficient path state for $r{=}1$ collapses to a tiny counter-based object.
%
For each active path $\TPath$, we only need to know whether it is alive, the number of remaining pivots on the path, and the fixed number of pivots still required to form a size-$s$ clique.
%
Hence the $r{=}1$ regime becomes an immutable-path instance of DPCM: the path transition is simple, and the support change of every surviving vertex is given by a closed-form difference.
%
The underlying path-state transition is described in Theorem~\ref{lem:path_vertex_remove}.
% \begin{lemma}[Vertex Removal]\label{lem:path_vertex_remove}
% Let $\TPath = \{V_{p}(\TPath), V_{h}(\TPath)\}$ be a clique path in the \cpi\ of a graph $G$, 
% where $V_{p}(\TPath)$ and $V_{h}(\TPath)$ denote the sets of pivot and hold vertices on $\TPath$, respectively. 
% %
% After removing a vertex $v$ from $G$, the corresponding update to $\TPath$ is as follows:
% % If a vertex $v$ is removed from $G$, the corresponding clique path $\TPath$ in the \cpi\ satisfies the following properties:
% \begin{enumerate}
%     \item If $v \in V_{h}(\TPath)$, the entire path $\TPath$ is removed from the \cpi.
%     \item If $v \in V_{p}(\TPath)$, the vertex $v$ is removed from $\TPath$, while the remaining path stays unchanged.
% \end{enumerate}
% \end{lemma}

\begin{theorem}[Vertex Removal]\label{lem:path_vertex_remove}
Let $\TPath$ be a clique path in the \cpi\ of a graph $G$, with hold and pivot
sets denoted by $V_h(\TPath)$ and $V_p(\TPath)$, respectively.
For any removed vertex $v \in \TPath$, the updated $\TPath'$ satisfies:

\begin{enumerate}[leftmargin=*]
    \item if $v \in V_h(\TPath)$, $\TPath' = \emptyset$;
    \item if $v \in V_p(\TPath)$, $\TPath' = \{(V_p(\TPath) \setminus \{v\},\, V_h(\TPath))\}$.
\end{enumerate}


% For all $s$-cliques encoded by $\TPath$ that do not contain $v$, they remain encoded by $\TPath'$.
\end{theorem}


\begin{proof}
By Lemma~\ref{lem:clique_count}, every $k$-clique encoded by a path
$\TPath{}$ must contain all hold vertices on that path, whereas pivot
vertices are optional. If the removed vertex $v$ lies in the hold set
$V_h(\TPath{})$, no encoded $k$-clique can survive, so the whole path
vanishes. If $v\in V_p(\TPath{})$, only those cliques that actually
include $v$ disappear; the cliques formed by the unchanged hold set and
the remaining pivots are still uniquely represented by the shortened
path, so we remove $v$ and keep the rest of $\TPath{}$ intact.
\end{proof}

% \begin{example}
% Let us consider the \cpi shown in Figure~\ref{fig:pruning_example}. For the \cpi shown in Figure~\ref{fig:pruning_example}, consider the clique path $\TPath{2}$ with
% $V_{p}(\TPath{2}) = \{v_4, v_5\}$ and $V_{h}(\TPath{2}) = \{v_1, v_6\}$. 
% If we remove a hold vertex $v_1$ from the underlying graph, then according to Lemma~\ref{lem:path_vertex_remove}, the entire path $\TPath{2}$ will be removed from the \cpi. On the other hand, if we remove a pivot vertex $v_4$ from the underlying graph, then according to Lemma~\ref{lem:path_vertex_remove}, we will remove $v_4$ from $\TPath{2}$, and the updated path will be $\TPath{2}'$ with $V_{p}(\TPath{2}') = \{v_5\}$ and $V_{h}(\TPath{2}') = \{v_1, v_6\}$.
% \end{example}

% \begin{algorithm}[t]
\small
\small
\caption{\textsc{UpdateIndex}(r=1)}
\label{alg:RemoveNode}
\KwIn{clique path $\TPath{} = \{V_{H},V_{P}\}$, vertices to remove $V_{\mathrm{rm}}$}
\KwOut{subpaths induced after deleting $V_{\mathrm{rm}}$}
\SetKwFunction{RemoveNode}{RemoveNode}
\SetKwFunction{CountPreVertices}{$s$-CliqueCount}
\SetKwProg{Fn}{Function}{:}{}


\lIf{$|V_{rm} \cap V_{H}(\TPath)| != 0$}{
    \Return
}
% output path $(V_{Piv}(\TPath) \setminus V, V_{Hold}(\TPath))$\;
\textbf{output path }$\{ V_{H}(\TPath),V_{P}(\TPath) \setminus V_{\mathrm{rm}}\}$\;
    %

\end{algorithm}

% \begin{algorithm}[t]
\small
% \caption{\textsc{RemoveNode} (r=1 path update)}
% \label{alg:RemoveNode}
% \KwIn{CPI $\tree$, clique path $\TPath$, vertex $v$, clique size $s$}
% \KwOut{Core value $\kappa_s(v)$ for every vertex $v$}
% \SetKwFunction{RemoveNode}{RemoveNode}
% \SetKwFunction{CountPreVertices}{$s$-CliqueCount}
% \SetKwProg{Fn}{Function}{:}{}


% \If{$v$ is hold in $\TPath$}{ 
%     remove $\TPath$ from $\tree$\;
% }
% \ElseIf{$v$ is pivot in $\TPath$}{
%     remove $v$ from $\TPath$\;
%     \If{$|\TPath| < s$}{
%         remove $\TPath$ from $\tree$\;
%     }
% }
% % }

\stitle{Batched Updates.}
% 我们可以很容易地将算法扩展到批量删除顶点的情况。给定一个顶点删除集合$V_{\mathrm{rm}}$和一个包含任何被删除顶点的路径$\TPath$，我们可以通过以下方式更新$\TPath$:
In each iteration of the peeling process, many $r$-cliques are removed. 
To avoid redundant computation, we adopt a batched strategy that merges repeated updates. 
For vertex removals, the algorithm can be easily extended to handle batches. 
Given a set of vertices to be removed $V_{\mathrm{rm}}$ and a path $\TPath$ containing any of them, we update $\TPath$ as follows:
\begin{enumerate}[leftmargin=*]
    \item $\TPath' = \emptyset$ if $|V_h(\TPath) \cap V_{\mathrm{rm}}| > 0$;
    \item $\TPath' = \{(V_p(\TPath) \setminus V_{\mathrm{rm}},\, V_h(\TPath))\}$ if $|V_h(\TPath) \cap V_{\mathrm{rm}}| = 0$.
\end{enumerate}
% 我们可以很轻松的根据上述规则更新$\TPath$。
We can easily update $\TPath$ according to the above rules.

\stitle{Closed-Form Contribution Delta.}
For $r{=}1$, the local support contribution of a path depends only on how many pivots remain after a peeling round.
%
This yields a fully analytical DPCM update rule.

\begin{theorem}[Closed-Form Support Delta for $r=1$]\label{thm:r1_delta}
Let $\TPath=(V_p,V_h)$ be an active path with $\mathit{rp}=|V_p|$ remaining pivots and $\mathit{np}=s-|V_h|$ pivots still needed to form a size-$s$ clique.
Suppose that a peeling round removes $d$ pivots from $\TPath$ and no hold vertex on $\TPath$ is removed.
If the path survives, i.e., $\mathit{np}\le \mathit{rp}-d$, then the support change of each surviving vertex $u$ contributed by $\TPath$ is
\[
\Delta_{\TPath}(u)=
\begin{cases}
\binom{\mathit{rp}}{\mathit{np}}-\binom{\mathit{rp}-d}{\mathit{np}}, & u\in V_h,\\[4pt]
\binom{\mathit{rp}-1}{\mathit{np}-1}-\binom{\mathit{rp}-d-1}{\mathit{np}-1}, & u\in V_p.
\end{cases}
\]
If a hold is removed or $\mathit{np}>\mathit{rp}-d$, then the path dies and its entire contribution is subtracted.
\end{theorem}

\begin{proof}
By Lemma~\ref{lem:clique_count}, a hold vertex belongs to every size-$s$ clique encoded by $\TPath$, so its path contribution equals the number of ways to choose the required pivots, namely $\binom{\mathit{rp}}{\mathit{np}}$ before the update and $\binom{\mathit{rp}-d}{\mathit{np}}$ after it.
%
For a surviving pivot vertex, one pivot is fixed by the vertex itself, so the corresponding counts become $\binom{\mathit{rp}-1}{\mathit{np}-1}$ and $\binom{\mathit{rp}-d-1}{\mathit{np}-1}$.
%
The support delta follows by subtraction.
\end{proof}

\stitle{Implication for DPCM.}
Theorem~\ref{thm:r1_delta} shows that the $r{=}1$ instantiation needs neither path splitting nor local clique reconstruction.
%
All exact support updates are reduced to constant-time binomial lookups on the maintained path counters.


% \end{algorithm}
% \stitle{Algorithm.} 
% Algorithm~\ref{alg:RemoveNode} implements the path update procedure described in Lemma~\ref{lem:path_vertex_remove},
% % 给定一个顶点删除集合$V_{\mathrm{rm}}$和一个包含任何被删除顶点的路径$\TPath$，我们根据上述规则更新$\TPath$。
% Given a set of vertices to be removed $V_{\mathrm{rm}}$ and a path $\TPath$ containing any of the removed vertices, we update $\TPath$ according to the rules above.


\stitle{Complexity.}
Let a clique path $\TPath = \{V_p(\TPath), V_h(\TPath)\}$.
For a single path update, the running time of removing a node from \cpi is $O\!\big(|V_h(\TPath)|+|V_p(\TPath)|\big)$ and the extra space is $O(1)$.
%
By Theorem~\ref{thm:r1_delta}, the differential support update on the path is fully analytical, so the overall update cost stays linear in the path length.




% \begin{table}[t]
% \centering
% \small
% \setlength{\tabcolsep}{10pt}
% \renewcommand{\arraystretch}{1.12}
% \begin{tabular}{c l}
% \toprule
% $r$ & Instantiated update operator (peeled object) \\
% \midrule
% $1$   & remove a vertex $v$ \\
% $2$   & remove an edge $e=\{u,v\}$ \\
% $\ge 3$ & remove an $r$-clique $R$ \\
% \bottomrule
% \end{tabular}
% \caption{A compact mapping from $r$ to the instantiated update operator.
% The maintained CPI state is identical across all $r$ (CPI paths with $(V_h,V_p)$, supports, and a peeling queue); only the peeled object differs.}
% \label{tab:mapping}
% \end{table}
